\section{Conclusion}
\label{sec:conclusion}

This survey examined existing research on how regional English accents may
affect Direct and Cascaded speech-to-text translation. Three main findings were
identified.

First, previous studies do not show that one speech-translation architecture is
always better. The results depend on factors such as language direction,
training resources, evaluation data, reference translations, and evaluation
metrics. Second, accent-related performance differences are well established
in ASR, including in model families related to modern speech-translation
systems. However, WER measures transcription error and does not directly
measure final translation quality. Third, research shows that ASR errors can
propagate through Cascaded systems and reduce translation quality. This makes
accent-related translation degradation possible, but existing studies do not
directly demonstrate it.

The closest studies still do not provide the full comparison required by the
research question. Some accent-level speech-translation benchmarks include
only Cascaded systems, while others do not combine matched regional-English
utterances, a complete Cascaded comparator, and accent-based final translation
results.

The datasets also have important limitations. Common Voice accent labels are
optional and self-reported rather than verified by expert phonetic analysis.
Different dataset versions and recording conditions may also affect the
results. Speakers or source sentences may appear more than once, and CoVoST 2
reference translations were created from written transcripts rather than
directly from audio. In addition, BLEU and chrF scores can change depending on
reference and processing choices, and small accent groups may make performance
differences uncertain.

Therefore, the research gap is not simply a lack of research on accent or
speech-translation robustness. The more specific gap is the absence of a
controlled, matched comparison where regional English accent is the
grouping variable and final translation quality is measured for both Direct
and Cascaded ST systems.

This supports the research question: How does regional English accent
affect the translation quality of Direct and Cascaded speech-to-text
translation systems?

A reliable evaluation requires both systems to use the same utterances and
reference translations, while speaker structure and accent groups should be
carefully controlled. Evaluation metrics should also be applied consistently,
and uncertainty should be measured using paired statistical analysis.

Such an evaluation can show whether system performance is associated with
accent group. However, any observed differences should not be explained only
by architecture, because the compared models may also differ in training data,
model capacity, objectives, and decoding methods.
